\section{The formal statements}

The companion Lean development separates its public definitions from
the proofs. The coverage predicate on a word $w$ over $\operatorname{Fin}K$
is the assertion that every list $p$ of length $K$ with distinct entries
has a prefix $u$ and a suffix $v$ with $w=upv$. Thus the final statements
refer directly to covering words, with no path or chain hypothesis.

The two statements in
\nolinkurl{Submissions/General/Challenge.lean}
are Theorem~\ref{thm:main} with the exact and quadratic return orders.
Their formulas include the initial rows $8\le K\le12$.
Their declaration names are \nolinkurl{exact_word_bound} and
\nolinkurl{quadratic_word_bound}, in the namespace
\nolinkurl{SuperpermutationResults.GeneralLower}.
The challenge imports only Mathlib. Its solution environment imports
proved declarations with the same names; the internal bridge between
the displayed formulas and the proof's formulas is checked by reflexivity.
The Lean formula uses the natural-number ceiling. The correction in
\eqref{eq:bound} is positive in every stated application, so this
agrees with the ordinary ceiling used here.

The challenge in \nolinkurl{Submissions/Asymptotics/Challenge.lean}
states thirteen consequences. The correspondence
with the paper is recorded below. Names in the second column are the
suffixes of declarations in
\nolinkurl{SuperpermutationPalomar.LowerAsymptotics}.

\begin{center}
\small
\begin{tabular}{@{}>{\raggedright\arraybackslash}p{0.37\textwidth}
                  >{\raggedright\arraybackslash}p{0.59\textwidth}@{}}
\toprule
Mathematical statement & Lean declaration\\
\midrule
The leading correction in Corollary~\ref{cor:comparison}
& \nolinkurl{exact_leading_limit}\newline
  \nolinkurl{quadratic_leading_limit}\\[3pt]
Corollary~\ref{cor:word-asymptotic}
& \nolinkurl{word_lower_bound_two_sub_epsilon}\\[3pt]
The three comparisons of $B$ with $H$
& \nolinkurl{exact_reference_limit_zero}\newline
  \nolinkurl{exact_reference_limit_one}\newline
  \nolinkurl{exact_reference_limit_two}\\[3pt]
The three comparisons of $Q$ with $H$
& \nolinkurl{quadratic_reference_limit_zero}\newline
  \nolinkurl{quadratic_reference_limit_one}\newline
  \nolinkurl{quadratic_reference_limit_two}\\[3pt]
The two limits for $Q-O$
& \nolinkurl{quadratic_gain_zero_limit}\newline
  \nolinkurl{quadratic_gain_one_limit}\\[3pt]
The two eventual lower estimates for $B-O$
& \nolinkurl{exact_gain_zero_asymptotic_lower}\newline
  \nolinkurl{exact_gain_one_asymptotic_lower}\\
\bottomrule
\end{tabular}
\end{center}

Here the suffixes zero, one and two refer to $n=K-1$ modulo three.
The last two declarations say that, for each positive epsilon, the
normalized difference is eventually at least the stated constant minus
that epsilon. They do not assert convergence. The finite strict-improvement
thresholds in Proposition~\ref{prop:gains} have the ordinary proof in
Appendix~\ref{app:numerical}; they are not separate claims in this
asymptotic challenge.

Every declaration just listed, and both general word bounds, has a proof
using only \nolinkurl{propext}, \nolinkurl{Classical.choice}, and
\nolinkurl{Quot.sound}. No native computation assumption is used in these
lower-bound families.
